\chapter{Proposed Method}

This chapter explains how a single kernel event becomes an alert. The processing
pipeline is shown in Figure~\ref{fig:pipeline}: an eBPF program captures the event in
the kernel, a C loader forwards it to the Python agent, the agent resolves its
Kubernetes identity, extracts features, applies the Random Forest and the rule engine,
computes a risk score and sends the result to the backend.

\begin{figure}[H]
\centering
\begin{tikzpicture}[node distance=0.55cm and 0.5cm, every node/.style={font=\small}]
  \node[kbox, minimum width=4.2cm] (k) {Kernel tracepoints\\ \scriptsize execve, openat, connect, ...};
  \node[kbox, below=of k, minimum width=4.2cm] (ebpf) {eBPF program (C, CO-RE)\\ \scriptsize in-kernel filtering};
  \node[kbox, below=of ebpf, minimum width=4.2cm] (rb) {BPF ring buffer $\rightarrow$ C loader\\ \scriptsize JSON lines, drop counters};
  \node[ubox, below=of rb, minimum width=4.2cm] (col) {Collector\\ \scriptsize normalise event};
  \node[ubox, below=of col, minimum width=4.2cm] (enr) {Kubernetes enrichment\\ \scriptsize cgroup $\rightarrow$ Pod / container / workload};
  \node[decision, below=0.5cm of enr] (ctx) {Pod context?};
  \node[dbox, right=1.2cm of ctx] (raw) {Store as raw event};
  \node[ubox, below=0.5cm of ctx, minimum width=4.2cm] (fe) {Feature extraction\\ \scriptsize 74 features + 60 s window};
  \node[mlbox, below left=0.6cm and -1.6cm of fe, minimum width=3.2cm] (rf) {Random Forest\\ \scriptsize $P(\mathrm{mal})$, category};
  \node[mlbox, below right=0.6cm and -1.6cm of fe, minimum width=3.2cm] (rules) {Rule engine\\ \scriptsize 19 rules};
  \node[mlbox, below=1.9cm of fe, minimum width=4.2cm] (risk) {Risk score 0--100\\ \scriptsize severity, MITRE technique};
  \node[dbox, below=of risk, minimum width=4.2cm] (be) {Backend: event / alert\\ \scriptsize de-duplication, MongoDB};

  \draw[arr] (k) -- (ebpf);
  \draw[arr] (ebpf) -- (rb);
  \draw[arr] (rb) -- (col);
  \draw[arr] (col) -- (enr);
  \draw[arr] (enr) -- (ctx);
  \draw[arr] (ctx) -- node[lbl, above]{no} (raw);
  \draw[arr] (ctx) -- node[lbl, right]{yes} (fe);
  \draw[arr] (fe) -- (rf);
  \draw[arr] (fe) -- (rules);
  \draw[arr] (rf) -- (risk);
  \draw[arr] (rules) -- (risk);
  \draw[arr] (risk) -- (be);
  \draw[arr] (raw) |- (be);

  \begin{scope}[on background layer]
    \node[draw, dashed, rounded corners, fill=orange!3, fit=(k)(rb), inner sep=6pt,
          label={[font=\scriptsize\itshape]left:kernel space}] {};
    \node[draw, dashed, rounded corners, fill=blue!2, fit=(col)(risk)(rf)(rules), inner sep=6pt,
          label={[font=\scriptsize\itshape]left:agent (user space)}] {};
  \end{scope}
\end{tikzpicture}
\caption{Processing pipeline of the proposed system}
\label{fig:pipeline}
\end{figure}

\section{Kernel-Level Event Capture with eBPF}

The sensor is an eBPF program (\texttt{bpf/ids.bpf.c}) compiled with clang and loaded
with libbpf using \emph{Compile Once -- Run Everywhere} (CO-RE). CO-RE relocates
kernel structure offsets at load time using the kernel's BTF information, so one binary
runs on any kernel version with BTF, without kernel headers on the node.

The program attaches to the system-call tracepoints in Table~\ref{tab:tracepoints}.
For every event it records the cgroup ID, PID, parent PID, user ID, process and parent
name, and the arguments relevant for that call. Several filters run \emph{inside} the
kernel to reduce overhead: \texttt{openat} events are only emitted for sensitive path
prefixes, \texttt{chmod} only when SUID/SGID or execute bits are set, and
\texttt{execve} is reported on syscall exit so that only \emph{successful} executions are
emitted. Process arguments are copied up to eight arguments of 127 bytes. Events are
written to a BPF ring buffer; if the buffer is full the event is counted in a per-CPU
drop counter, which the agent reports so that data loss is never silent.

\begin{table}[H]
\centering
\caption{Monitored system-call tracepoints}
\label{tab:tracepoints}
\footnotesize
\begin{tabularx}{\textwidth}{l >{\raggedright\arraybackslash}p{5.2cm} >{\raggedright\arraybackslash}X}
\toprule
\textbf{Event type} & \textbf{Tracepoint(s)} & \textbf{Data captured} \\
\midrule
exec    & sys\_enter / sys\_exit\_execve & binary path, argv (8 $\times$ 127 B) \\
open    & sys\_enter\_openat & path, flags (sensitive prefixes only) \\
connect & sys\_enter\_connect & family, port, IPv4/IPv6 address, UNIX path \\
ptrace  & sys\_enter\_ptrace & request, target PID \\
setuid  & sys\_enter\_setuid, sys\_enter\_setresuid & target user ID \\
mount   & sys\_enter\_mount & source, target, filesystem type \\
unshare & sys\_enter\_unshare & namespace flags \\
setns   & sys\_enter\_setns & namespace type \\
chmod   & sys\_enter\_chmod, fchmodat, fchmodat2 & path, mode \\
module  & sys\_enter\_init\_module, finit\_module & module load \\
memfd   & sys\_enter\_memfd\_create & memfd name \\
\midrule
\multicolumn{3}{l}{\textbf{16 tracepoints in total}} \\
\bottomrule
\end{tabularx}
\end{table}

A small C program (\texttt{loader/ids\_loader.c}) loads and attaches the eBPF program,
polls the ring buffer and prints each event as one JSON line on standard output, where
the Python collector reads it.

\section{Kubernetes Enrichment}

Kernel events identify a container only through its \textbf{cgroup ID}. The enrichment
module resolves the cgroup ID to the cgroup path under \texttt{/sys/fs/cgroup}. Regular
expressions extract the \textbf{Pod UID} and the \textbf{container ID} from the path
(for the systemd and cgroupfs layouts used by containerd and Docker). A cache of Pod
objects on the node, refreshed from the Kubernetes API every 15 seconds, maps the Pod
UID to the Pod name, namespace, container name, image, service account and node. The
owning \textbf{workload} (Deployment, StatefulSet, DaemonSet, Job) is derived from owner
references, for example a ReplicaSet \texttt{shop-web-7b955dcf77} belongs to
Deployment \texttt{shop-web}. Events from processes outside any pod (the node itself)
keep an empty namespace and are stored as \emph{raw} events.

\section{Feature Extraction}

The feature module (\texttt{ids\_agent/features.py}) converts an event and its
behaviour window into a vector of \textbf{74 numeric features}. The same module is
imported by the training script and by the live agent, which removes the
train/serve mismatch of the v0 prototype. Table~\ref{tab:features} summarises the
feature groups.

\begin{table}[H]
\centering
\caption{Feature groups (74 features)}
\label{tab:features}
\small
\begin{tabularx}{\textwidth}{>{\raggedright\arraybackslash}p{3.3cm} c >{\raggedright\arraybackslash}X}
\toprule
\textbf{Group} & \textbf{\#} & \textbf{Examples} \\
\midrule
Event type (one-hot) & 11 & ev\_exec, ev\_open, ev\_connect, ev\_mount, ev\_memfd \\
Identity & 2 & uid\_root, ns\_system (platform namespace) \\
Command line & 9 & length, number of arguments, special characters, Shannon entropy, base64 blob, IP literal, URL, \texttt{/dev/tcp} \\
Binary class & 10 & shell, network tool, interpreter, recon tool, container tool, escape tool, package manager, runs from /tmp, hidden, fileless \\
Process lineage & 7 & K8s component, auth binary, debugger, parent is shell / server / interpreter / K8s component \\
Argument patterns & 9 & sensitive path, download, download-and-execute, write to /tmp, SUID, reverse shell, miner, log cleanup, recon \\
File & 2 & sensitivity level of path (0--2), write access \\
Network & 7 & UNIX socket, common / suspicious / high port, private, loopback, link-local destination \\
Privilege & 9 & ptrace attach, setuid to root, user / mount namespace unshare, SUID / exec chmod, benign memfd name, block-device mount, host-path mount \\
Behaviour window & 8 & executions, distinct binaries, connections, distinct destinations, sensitive opens, shells, new binary, shell ratio \\
\bottomrule
\end{tabularx}
\end{table}

\textbf{Command-line entropy.} For a command line $s$ with character frequencies $p_c$,
the Shannon entropy
\begin{equation}
H(s) = -\sum_{c} p_c \log_2 p_c
\end{equation}
is high for encoded or random payloads and low for ordinary commands.

\textbf{Behaviour window.} For each container the agent keeps a sliding window of the
last $W = 60$ seconds of events. The window counts are log-transformed,
$x' = \log(1 + x)$, so that bursts of thousands of events do not dominate. A binary is
marked as \emph{new} if it has not been executed in the same workload before, once a
warm-up period of five minutes has passed; this captures the classic sign of an
intrusion: a container that suddenly runs programs it never ran before.

\section{Random Forest Classifier}

A Random Forest~\cite{breiman2001} is an ensemble of $T$ decision trees, each trained on
a bootstrap sample of the training data with a random subset of features considered at
every split. For an input vector $\mathbf{x}$, every tree $t$ outputs class
probabilities $p_t(c \mid \mathbf{x})$ and the forest averages them:
\begin{equation}
P(c \mid \mathbf{x}) = \frac{1}{T} \sum_{t=1}^{T} p_t(c \mid \mathbf{x}).
\end{equation}
The model is trained on \textbf{ten classes} (benign and nine attack categories). Two
quantities are derived from the output:
\begin{align}
P(\mathrm{mal} \mid \mathbf{x}) &= 1 - P(\mathrm{benign} \mid \mathbf{x}), \\
\hat{c}(\mathbf{x}) &= \arg\max_{c \neq \mathrm{benign}} P(c \mid \mathbf{x}).
\end{align}
An event is an ML detection if $P(\mathrm{mal}) \geq \tau$ with $\tau = 0.5$; in
platform namespaces (for example \texttt{kube-system}) without a matching rule, the
threshold is raised to 0.75 because platform components legitimately perform many
sensitive operations. The predicted attack category $\hat{c}$ selects the MITRE
technique of an ML-only alert.

\textbf{Hyper-parameters:} 150 trees, maximum depth 22, minimum one sample per leaf,
$\sqrt{74}$ features per split, no class weighting, random state 42. Balanced class
weights were tested and rejected: they raised detection to 100\% but increased the
false-positive rate to 4.6\%, because up-weighting nine small attack classes inflates
$P(\mathrm{mal})$.

\section{Rule Engine}

Rules provide high-precision detection of well-known techniques and an explanation
that operators understand. Each rule receives the event and its feature vector and
returns true or false; a failing rule is skipped and never stops the sensor. Platform
components (Kubernetes namespaces such as \texttt{kube-system}, container runtimes,
CoreDNS) are exempted where their normal behaviour would otherwise match.

\begin{longtable}{p{1.0cm} p{6.6cm} p{1.5cm} p{2.2cm}}
\caption{Detection rules}\label{tab:rules}\\
\toprule
\textbf{ID} & \textbf{Rule} & \textbf{Severity} & \textbf{MITRE} \\
\midrule
\endfirsthead
\toprule
\textbf{ID} & \textbf{Rule} & \textbf{Severity} & \textbf{MITRE} \\
\midrule
\endhead
R001 & Reverse-shell command line & High & T1059.004 \\
R002 & Shell / interpreter connecting to a suspicious port & High & T1571 \\
R003 & Service-account token read by a shell tool & High & T1528 \\
R004 & Password database read & High & T1003.008 \\
R005 & Crypto-miner activity (arguments or pool ports) & High & T1496 \\
R006 & Download piped into a shell & High & T1105 \\
R007 & Binary executed from a temp / hidden location & Medium & T1204.002 \\
R008 & Kernel module loaded from a pod & High & T1547.006 \\
R009 & Process injection via ptrace & Medium & T1055.008 \\
R010 & SUID/SGID bit set on a binary & High & T1548.001 \\
R011 & Web / server process spawned a shell or recon tool & Medium & T1190 \\
R012 & Fileless execution (memfd / fd exec) & High & T1620 \\
R013 & Log or shell-history tampering & Medium & T1070 \\
R014 & Container-escape technique (nsenter to PID 1, cgroup / block-device mount, user-namespace unshare, runtime socket, core\_pattern) & High & T1611 \\
R015 & Cloud metadata endpoint access & Medium & T1552.005 \\
R016 & Persistence / authentication file modified & High & T1098 \\
R017 & Kubelet / etcd API accessed from a workload & High & T1210 \\
R018 & Network scanning tool & Medium & T1046 \\
R019 & Privilege escalation to root by a non-auth binary & High & T1068 \\
\bottomrule
\end{longtable}

\section{Risk Scoring and Severity}

ML and rules are combined into a single risk score. Let $H$ be the set of matched
rules, $s(r)$ the score of rule $r$ ($\mathrm{High}=90$, $\mathrm{Medium}=70$,
$\mathrm{Low}=45$) and $m = \mathbb{1}[P(\mathrm{mal}) \geq \tau]$ the ML decision. Then
\begin{equation}
S_0 = \max\Big( m \cdot 100\,P(\mathrm{mal}),\; \max_{r \in H} s(r) \Big)
\end{equation}
and the final score adds context bonuses and penalties and is clipped to $[0, 100]$:
\begin{equation}
S = \mathrm{clip}\big(S_0 + 5\cdot\mathbb{1}[H \neq \emptyset \wedge m]
      + 3\cdot\mathbb{1}[\mathrm{uid}=0] - 10\cdot\mathbb{1}[\mathrm{ns\_system} \wedge H = \emptyset],\; 0,\; 100\big).
\end{equation}
The bonus of 5 rewards agreement of both detectors, running as root adds 3, and an
ML-only verdict in a platform namespace loses 10. The event is \emph{suspicious} if
$H \neq \emptyset$ or $m = 1$. Severity is assigned as
\begin{equation}
\mathrm{severity}(S) =
\begin{cases}
\mathrm{High}   & S \geq 80,\\
\mathrm{Medium} & 60 \leq S < 80,\\
\mathrm{Low}    & S < 60.
\end{cases}
\end{equation}
The alert's category comes from the highest-severity matched rule, or from the ML
prediction $\hat{c}$ if no rule matched. Each alert lists the matched rules, the ML
probability, the MITRE techniques and the full Kubernetes context.

\section{Alert De-duplication}

An attacker who repeats a command, or a loop that runs every second, would otherwise
create hundreds of identical alerts. The agent computes a de-duplication key as a SHA-1
hash of namespace, workload, container, event type, process name, category and the
command line (or path or destination). The backend merges a new alert into an
unresolved alert with the same key seen within the last \textbf{10 minutes}, increments its
\emph{count} and updates its last-seen time and maximum risk score.
